\cotefile{157-4}
\scanpage{66}
\mspage{59}
% archive sheet number: 65
\uline{Théorème-scholie} : Soit $(C,\mathrm{TC},F,\mathrm{TF})$ une \unsure{quillénisation} uniforme gauche de $W$ sur $\mathrm{Diag}\subset\mathrm{Cat}$.

1) Soit $u:I'\unsure{-|\to}I$ une flèche de $\mathrm{Diag}$. On \deleted{\unsure{supposera que}} \unsure{$\mathrm{Diag}$ accordé|$\mathrm{Diag}$ associé} \unsure{à|:} $M$ de façon que pour $u$, flèche de $\mathrm{Diag}$, les $\varprojlim$ dans $M$ \deleted{\unsure{requises}} (de type $I'_{/i}$) requises pour assurer l'existence de
\[
u_*:M(I')\longrightarrow M(I)
\]
existent dans $M$. (C'est OK si $I^{\deleted{\unsure{'|*}}},I'$ finis, $M$ étant \unsure{supposé} stable par $\varprojlim$ finies de toute façon.) On sait (prop. 1 page 45) que $u_*$ induit
\[
u^{f}_{*\deleted{\unsure{f}}}:M(I')_f\longrightarrow M(I)_f
\tag*{(*)}
\]
On dit que $u$ est \deleted{\unsure{$W$-}}\uline{admissible} \inserted{(à gauche)} pour la quillénisation uniforme gauche donnée \inserted{de $W$ sur $\mathrm{Diag}$} si $u^f_*$ est compatible avec les localiseurs $W(I')_f$, $W(I)_f$ i.e. si $f':X'\unsure{-|\to}Y'$ dans $M(I')$, $X',Y'$ \deleted{$\unsure{\in}$} fibrants, $f'\in W(I')$ implique $u_*\unsure{(f)|(f')}\in W(I)$.

On dit que $Z$ dans $\mathrm{Diag}$ est admissible (pour la quillénisation uniforme gauche donnée de $W$ sur $\mathrm{Diag}$\ldots) si le foncteur $Z\unsure{-|\to}e$ est admissible. Ceci posé,

a) Pour que $u$ soit admissible, il suffit que les $I'_{/s}$ pour $s\in\mathrm{Ob}\,I$, le soient.
\margin{NB il y a une réciproque, quand $I$ est \deleted{$\unsure{I|s\backslash}$} ordonné, donc les \deleted{$\unsure{I|s\backslash}$} $I'_{/s}\unsure{-|\to}I'$ des immers. ouvertes.}
